% 第1章 引言
\section{引言：缝隙在哪里}\label{sec:intro}

先算一笔小账。一条样本内年化 $20\%$ 的多空信号，年换手 $100$ 倍，单边成本 $10$ 个基点——光交易成本一年就收走 $20$ 个百分点。信号还在，钱没了。这笔账谁都算得出来，但它容易引出两个错误的推论：一是“信号是假的”，二是“交易是徒劳的”。更接近事实的读法是第三种：\textbf{可预测性与可交易性之间横着一段距离，这段距离的性质，决定了量化研究里几乎每一条结论的成色。}

这段距离由三段落差叠成。第一段是统计的：候选池里挑出回测最好的那条，“最好”本身意味着向上偏差，要校正它，需要多重检验的纪律\cite{white2000,harveylz2016}。第二段是摩擦的：佣金、税费、冲击成本、做空可行性，每一项都要从收益里扣，而且彼此不是简单相加。第三段最少被形式化，却是本文的主题——\textbf{观测的落差}：一个信号能不能变成钱，取决于你在什么时间尺度上观测它，以及在什么参照系里度量它。

后两件事的例子随手就有。时间尺度这一侧：同一条反转信号，日频调仓时交易成本比信号幅度还大；把调仓周期放宽到数日，成本被按比例摊薄，而预测力尚未完全衰减，净值就从亏损翻成接近盈亏平衡——本文第 \ref{sec:empirical} 章会给出这条完整曲线。参照系这一侧：一个组合对等权基准的 $\alpha$ 显著为正，把基准换成市值加权指数或加入动量因子，同一段收益的 $\alpha$ 可能缩水大半。前者说明“观测频率”不是实现细节，它是缝隙存在性的一个自变量；后者说明“缝隙大小”不是一个绝对数，它相对于比较工具而定义。

这两件事在实务里是常识，在文献里也各有线索。异质市场假说早就提出，市场由不同时间尺度的参与者构成，时间尺度本身是市场结构的一部分\cite{muller1993}；基准批评\cite{roll1977}与联合假设问题\cite{fama1970}则早就指出，$\alpha$ 检验无法与定价模型的假设分离。但这些线索分散在不同传统里：时间尺度一侧主要停在实证与描述层面，参照系一侧主要停在资产定价的检验语境里，二者很少被合成同一个存在性命题。本文要做的事，就是用最小模型把它们收拢：给“可交易缝隙”一个定义，给它的存在性两个结构性命题——一个关于时间尺度，一个关于参照系——以及一条必要条件（信息幅度与成本的赛跑）；然后把实务中常用的筛选流程（本文称之为“三道门”，见第 \ref{sec:criteria} 章）与统计门槛解释为这些条件的操作版本；最后用真实数据把命题的可检验推论跑出来。

定位需要说清：本文不是文献综述，尽管吸收了若干文献结论；也不是实现教程，尽管所附实证可以按脚本与数据口径复现。它是一篇带实证的理论文章，核心是两个可证的命题，其余部分都在为这两个命题做铺垫与收尾。若读者只带走一句话，应该是这句：\textbf{“这条信号有没有缝隙”是一个不完整的问句，完整的问句必须带上尺度与参照系——“在什么观测频率上、相对于什么基准、扣除什么成本之后。”}

\subsection{用词说明}

“缝隙”与“可交易缝隙”是本文的工作术语，不对应文献中任何一个既有概念的完整含义。之所以自造而不直接借用现有词汇，是因为既有概念各自缺一块：\emph{mispricing}（定价偏差）不要求偏差能被提取；\emph{limits to arbitrage}（套利限制）解释偏差为何不消失，但不给出“值不值得动手”的判据；\emph{alpha} 相对基准定义，却通常不携带多重检验与摩擦的限定。本文要的语言需要同时容纳这三块。表 \ref{tab:terms} 给出对齐关系，避免读者把本文结论误配到相邻概念上。

\begin{table}[htbp]
\centering
\caption{本文术语与既有概念的对齐}
\label{tab:terms}
\begin{tabular}{@{}>{\raggedright\arraybackslash}p{2.6cm}>{\raggedright\arraybackslash}p{4.2cm}>{\raggedright\arraybackslash}p{7.6cm}@{}}
\toprule
本文术语 & 最接近的既有概念 & 差异与取舍 \\
\midrule
缝隙（gap） & mispricing、anomaly、alpha & 本文特指“扣除摩擦后仍为正、且经得起选择偏差校正”的残余收益源；三者各缺一项限定 \\[0.3em]
可交易缝隙 & tradable alpha、可实施空间 & 强调三个附加条件：多重检验校正后显著、样本外与跨样本不反号、成本按真实换手逐笔收取 \\[0.3em]
观测装置 & 无单一对应；邻近：异质市场假说的“参与者尺度”、Roll 基准批评中的“基准不可观测” & 本文把它形式化为二元组（时间尺度 $\mathbb{T}$，参照系 $\mathcal{F}$），使“相对于装置的存在性”可以被证明 \\[0.3em]
观测缝隙 & 无既有对应 & 本文新造的判据概念：一个缝隙声明在给定观测装置下是否成立 \\
\bottomrule
\end{tabular}
\end{table}

\subsection{阅读导览}

第 \ref{sec:economics} 章铺地基：有效市场之争留下的真正结论、信息租金为何耗散、以及低信噪比对方法论的三条压迫。第 \ref{sec:definition} 章给出可交易缝隙的定义与必要条件（信息幅度必须超过成本），并把判据改写成带估计误差的置信下界形式。第 \ref{sec:measurement} 章与第 \ref{sec:reference} 章是全文的两个核心：前者把时间尺度做成“尺度窗口”命题，后者把参照系做成“分解公式”。第 \ref{sec:criteria} 章把前两章的结论合成为四道核验，并解释它与实务中“三道门”的映射。第 \ref{sec:empirical} 章是实证：八年 A 股日线上的尺度扫描与参照系分解。第 \ref{sec:limits} 章反过来自查：判据在什么情形下会骗人。第 \ref{sec:conclusion} 章收束。